\documentclass[UTF8,a4paper,zihao=-4,fontset=none]{ctexart}
\newcommand{\ReportHeadLeft}{因子赛道 ｜ AI 工具使用说明}
\newcommand{\ReportHeadRight}{2026 年 11 月}
% 决赛策略研究报告 公共导言（XeLaTeX + ctexart）
% 原则：层级清楚（章、节两级为主，最多到小节）；正文成段；不用彩色框、侧栏、底纹、小字注释。
% 图表标题在上，写成「图 1　标题（数据范围）」，数据来源写进标题括号或正文，不另起小字行。
%
% 用法：
%   \documentclass[UTF8,a4paper,zihao=-4,fontset=none]{ctexart}
%   % 决赛策略研究报告 公共导言（XeLaTeX + ctexart）
% 原则：层级清楚（章、节两级为主，最多到小节）；正文成段；不用彩色框、侧栏、底纹、小字注释。
% 图表标题在上，写成「图 1　标题（数据范围）」，数据来源写进标题括号或正文，不另起小字行。
%
% 用法：
%   \documentclass[UTF8,a4paper,zihao=-4,fontset=none]{ctexart}
%   % 决赛策略研究报告 公共导言（XeLaTeX + ctexart）
% 原则：层级清楚（章、节两级为主，最多到小节）；正文成段；不用彩色框、侧栏、底纹、小字注释。
% 图表标题在上，写成「图 1　标题（数据范围）」，数据来源写进标题括号或正文，不另起小字行。
%
% 用法：
%   \documentclass[UTF8,a4paper,zihao=-4,fontset=none]{ctexart}
%   \input{../_style/report_preamble.tex}
%   \title{…} \subtitle{…} \team{…} \reportdate{…}
%   \begin{document} \makereporttitle ...

\setCJKmainfont{Songti SC}[BoldFont={Songti SC Bold}]
\def\yhpath{/Applications/Microsoft PowerPoint.app/Contents/Resources/DFonts/}
\setCJKsansfont{msyh.ttc}[Path=\yhpath,BoldFont=msyhbd.ttc]
\setCJKmonofont{msyh.ttc}[Path=\yhpath]
\setmainfont{Times New Roman}
\setsansfont{msyh.ttc}[Path=\yhpath,BoldFont=msyhbd.ttc]
\setmonofont{Menlo}[Scale=0.88]

\usepackage{newunicodechar}
\newunicodechar{→}{\ensuremath{\rightarrow}}
\newunicodechar{≈}{\ensuremath{\approx}}
\newunicodechar{≤}{\ensuremath{\leq}}
\newunicodechar{≥}{\ensuremath{\geq}}
\newunicodechar{−}{\ensuremath{-}}
\newunicodechar{×}{\ensuremath{\times}}
\newunicodechar{÷}{\ensuremath{\div}}
\newunicodechar{±}{\ensuremath{\pm}}
\newunicodechar{β}{\ensuremath{\beta}}
\newunicodechar{λ}{\ensuremath{\lambda}}
\newunicodechar{σ}{\ensuremath{\sigma}}

\usepackage{xcolor}
\definecolor{rnavy}{HTML}{1B365D}
\definecolor{rgrey}{HTML}{7A828C}
\definecolor{rline}{HTML}{BFC5CC}
\definecolor{raccent}{HTML}{C0392B}

\usepackage[a4paper,top=2.6cm,bottom=2.6cm,left=2.7cm,right=2.7cm,headheight=14pt]{geometry}
\linespread{1.45}
\setlength{\parskip}{0.35em}

% 标题层级：章（\section）、节（\subsection）、小节（\subsubsection，少用）
\ctexset{
  section={format=\sffamily\Large\bfseries\color{rnavy}\raggedright, name={,}, number=\arabic{section},
           aftername=\enspace, beforeskip=1.6em, afterskip=0.8em,
           numberformat={\color{white}\colorbox{rnavy}{\makebox[1.1em]{\arabic{section}}}}},
  subsection={format=\sffamily\large\bfseries\color{rnavy}\raggedright, beforeskip=1.2em, afterskip=0.5em},
  subsubsection={format=\sffamily\normalsize\bfseries\raggedright, beforeskip=0.9em, afterskip=0.3em},
}

% 页眉：左栏目、右日期，下方细色线；页脚：页码
\usepackage{fancyhdr}
\providecommand{\ReportHeadLeft}{策略研究报告}
\providecommand{\ReportHeadRight}{2026 年 11 月}
\pagestyle{fancy}
\fancyhf{}
\fancyhead[L]{\sffamily\small\color{rnavy}\ReportHeadLeft}
\fancyhead[R]{\sffamily\small\color{rnavy}\ReportHeadRight}
\renewcommand{\headrule}{{\color{rnavy}\hrule height 0.8pt}}
\fancyfoot[C]{\sffamily\small\thepage}
\fancypagestyle{plain}{\fancyhf{}\renewcommand{\headrulewidth}{0pt}\fancyfoot[C]{\sffamily\small\thepage}}

% 图表：标题在上，「图 1　标题」
\usepackage{graphicx,float}
\usepackage[font=small,labelfont={sf,bf,color=rnavy},textfont={sf,bf,color=rnavy},labelsep=colon,justification=raggedright,singlelinecheck=false,skip=5pt]{caption}
\DeclareCaptionLabelSeparator{zhcolon}{：}
\captionsetup{labelsep=zhcolon}
\captionsetup[figure]{name=图表,position=top}
\captionsetup[table]{name=图表,position=top}
\makeatletter\let\c@table\c@figure\makeatother
% 图表下方一行资料来源（唯一允许的小字）
\newcommand{\source}[1]{\par\vspace{2pt}{\raggedright\sffamily\footnotesize\color{rgrey}资料来源：#1\par}}

\usepackage{booktabs,tabularx,array,multirow}
\renewcommand{\arraystretch}{1.3}
\newcolumntype{L}{>{\raggedright\arraybackslash}X}
\newcolumntype{C}{>{\centering\arraybackslash}X}

\usepackage{amsmath,amssymb,bm}
\usepackage{enumitem}
\setlist{nosep,leftmargin=2em,itemsep=0.2em}

\usepackage{tikz}
\usetikzlibrary{arrows.meta,positioning,calc,fit,decorations.pathreplacing}
\tikzset{
  node/.style = {draw=rnavy, line width=0.6pt, rounded corners=1.5pt, align=flush center,
                 inner xsep=6pt, inner ysep=4pt, font=\sffamily\small},
  key/.style  = {node, draw=raccent, line width=0.8pt},
  soft/.style = {node, draw=rline},
  flow/.style = {-{Stealth[length=2mm]}, line width=0.6pt, draw=rnavy},
  note/.style = {font=\sffamily\small, text=rgrey, align=flush center},
}

% 首页研究要点：每条以一句加粗结论开头
\newenvironment{keypoints}{%
  \par\noindent{\sffamily\large\bfseries\color{rnavy}研究要点}\par\vspace{0.3em}
  {\color{rnavy}\hrule height 0.6pt}\vspace{0.6em}
  \begin{itemize}[leftmargin=1.2em,itemsep=0.6em,label={\color{rnavy}\rule[0.3ex]{0.45em}{0.45em}}]}
  {\end{itemize}\vspace{0.4em}{\color{rline}\hrule height 0.4pt}}
\newcommand{\kp}[1]{{\sffamily\bfseries\color{rnavy}#1}}

\usepackage[hidelinks]{hyperref}
\usepackage{bookmark}

% 封面
\makeatletter
\newcommand{\subtitle}[1]{\gdef\@subtitle{#1}}
\newcommand{\team}[1]{\gdef\@team{#1}}
\newcommand{\reportdate}[1]{\gdef\@reportdate{#1}}
\providecommand{\@subtitle}{}\providecommand{\@team}{}\providecommand{\@reportdate}{}
\newcommand{\makereporttitle}{%
  \thispagestyle{fancy}
  \vspace*{0.4em}
  {\raggedright\sffamily\color{rnavy}\fontsize{20pt}{27pt}\selectfont\bfseries\@title\par}
  \vspace{0.5em}
  {\raggedright\sffamily\large\color{rnavy}\@subtitle\par}
  \vspace{0.3em}
  {\raggedright\sffamily\small\color{rgrey}\@team\qquad\@reportdate\par}
  \vspace{1.2em}}
\makeatother

%   \title{…} \subtitle{…} \team{…} \reportdate{…}
%   \begin{document} \makereporttitle ...

\setCJKmainfont{Songti SC}[BoldFont={Songti SC Bold}]
\def\yhpath{/Applications/Microsoft PowerPoint.app/Contents/Resources/DFonts/}
\setCJKsansfont{msyh.ttc}[Path=\yhpath,BoldFont=msyhbd.ttc]
\setCJKmonofont{msyh.ttc}[Path=\yhpath]
\setmainfont{Times New Roman}
\setsansfont{msyh.ttc}[Path=\yhpath,BoldFont=msyhbd.ttc]
\setmonofont{Menlo}[Scale=0.88]

\usepackage{newunicodechar}
\newunicodechar{→}{\ensuremath{\rightarrow}}
\newunicodechar{≈}{\ensuremath{\approx}}
\newunicodechar{≤}{\ensuremath{\leq}}
\newunicodechar{≥}{\ensuremath{\geq}}
\newunicodechar{−}{\ensuremath{-}}
\newunicodechar{×}{\ensuremath{\times}}
\newunicodechar{÷}{\ensuremath{\div}}
\newunicodechar{±}{\ensuremath{\pm}}
\newunicodechar{β}{\ensuremath{\beta}}
\newunicodechar{λ}{\ensuremath{\lambda}}
\newunicodechar{σ}{\ensuremath{\sigma}}

\usepackage{xcolor}
\definecolor{rnavy}{HTML}{1B365D}
\definecolor{rgrey}{HTML}{7A828C}
\definecolor{rline}{HTML}{BFC5CC}
\definecolor{raccent}{HTML}{C0392B}

\usepackage[a4paper,top=2.6cm,bottom=2.6cm,left=2.7cm,right=2.7cm,headheight=14pt]{geometry}
\linespread{1.45}
\setlength{\parskip}{0.35em}

% 标题层级：章（\section）、节（\subsection）、小节（\subsubsection，少用）
\ctexset{
  section={format=\sffamily\Large\bfseries\color{rnavy}\raggedright, name={,}, number=\arabic{section},
           aftername=\enspace, beforeskip=1.6em, afterskip=0.8em,
           numberformat={\color{white}\colorbox{rnavy}{\makebox[1.1em]{\arabic{section}}}}},
  subsection={format=\sffamily\large\bfseries\color{rnavy}\raggedright, beforeskip=1.2em, afterskip=0.5em},
  subsubsection={format=\sffamily\normalsize\bfseries\raggedright, beforeskip=0.9em, afterskip=0.3em},
}

% 页眉：左栏目、右日期，下方细色线；页脚：页码
\usepackage{fancyhdr}
\providecommand{\ReportHeadLeft}{策略研究报告}
\providecommand{\ReportHeadRight}{2026 年 11 月}
\pagestyle{fancy}
\fancyhf{}
\fancyhead[L]{\sffamily\small\color{rnavy}\ReportHeadLeft}
\fancyhead[R]{\sffamily\small\color{rnavy}\ReportHeadRight}
\renewcommand{\headrule}{{\color{rnavy}\hrule height 0.8pt}}
\fancyfoot[C]{\sffamily\small\thepage}
\fancypagestyle{plain}{\fancyhf{}\renewcommand{\headrulewidth}{0pt}\fancyfoot[C]{\sffamily\small\thepage}}

% 图表：标题在上，「图 1　标题」
\usepackage{graphicx,float}
\usepackage[font=small,labelfont={sf,bf,color=rnavy},textfont={sf,bf,color=rnavy},labelsep=colon,justification=raggedright,singlelinecheck=false,skip=5pt]{caption}
\DeclareCaptionLabelSeparator{zhcolon}{：}
\captionsetup{labelsep=zhcolon}
\captionsetup[figure]{name=图表,position=top}
\captionsetup[table]{name=图表,position=top}
\makeatletter\let\c@table\c@figure\makeatother
% 图表下方一行资料来源（唯一允许的小字）
\newcommand{\source}[1]{\par\vspace{2pt}{\raggedright\sffamily\footnotesize\color{rgrey}资料来源：#1\par}}

\usepackage{booktabs,tabularx,array,multirow}
\renewcommand{\arraystretch}{1.3}
\newcolumntype{L}{>{\raggedright\arraybackslash}X}
\newcolumntype{C}{>{\centering\arraybackslash}X}

\usepackage{amsmath,amssymb,bm}
\usepackage{enumitem}
\setlist{nosep,leftmargin=2em,itemsep=0.2em}

\usepackage{tikz}
\usetikzlibrary{arrows.meta,positioning,calc,fit,decorations.pathreplacing}
\tikzset{
  node/.style = {draw=rnavy, line width=0.6pt, rounded corners=1.5pt, align=flush center,
                 inner xsep=6pt, inner ysep=4pt, font=\sffamily\small},
  key/.style  = {node, draw=raccent, line width=0.8pt},
  soft/.style = {node, draw=rline},
  flow/.style = {-{Stealth[length=2mm]}, line width=0.6pt, draw=rnavy},
  note/.style = {font=\sffamily\small, text=rgrey, align=flush center},
}

% 首页研究要点：每条以一句加粗结论开头
\newenvironment{keypoints}{%
  \par\noindent{\sffamily\large\bfseries\color{rnavy}研究要点}\par\vspace{0.3em}
  {\color{rnavy}\hrule height 0.6pt}\vspace{0.6em}
  \begin{itemize}[leftmargin=1.2em,itemsep=0.6em,label={\color{rnavy}\rule[0.3ex]{0.45em}{0.45em}}]}
  {\end{itemize}\vspace{0.4em}{\color{rline}\hrule height 0.4pt}}
\newcommand{\kp}[1]{{\sffamily\bfseries\color{rnavy}#1}}

\usepackage[hidelinks]{hyperref}
\usepackage{bookmark}

% 封面
\makeatletter
\newcommand{\subtitle}[1]{\gdef\@subtitle{#1}}
\newcommand{\team}[1]{\gdef\@team{#1}}
\newcommand{\reportdate}[1]{\gdef\@reportdate{#1}}
\providecommand{\@subtitle}{}\providecommand{\@team}{}\providecommand{\@reportdate}{}
\newcommand{\makereporttitle}{%
  \thispagestyle{fancy}
  \vspace*{0.4em}
  {\raggedright\sffamily\color{rnavy}\fontsize{20pt}{27pt}\selectfont\bfseries\@title\par}
  \vspace{0.5em}
  {\raggedright\sffamily\large\color{rnavy}\@subtitle\par}
  \vspace{0.3em}
  {\raggedright\sffamily\small\color{rgrey}\@team\qquad\@reportdate\par}
  \vspace{1.2em}}
\makeatother

%   \title{…} \subtitle{…} \team{…} \reportdate{…}
%   \begin{document} \makereporttitle ...

\setCJKmainfont{Songti SC}[BoldFont={Songti SC Bold}]
\def\yhpath{/Applications/Microsoft PowerPoint.app/Contents/Resources/DFonts/}
\setCJKsansfont{msyh.ttc}[Path=\yhpath,BoldFont=msyhbd.ttc]
\setCJKmonofont{msyh.ttc}[Path=\yhpath]
\setmainfont{Times New Roman}
\setsansfont{msyh.ttc}[Path=\yhpath,BoldFont=msyhbd.ttc]
\setmonofont{Menlo}[Scale=0.88]

\usepackage{newunicodechar}
\newunicodechar{→}{\ensuremath{\rightarrow}}
\newunicodechar{≈}{\ensuremath{\approx}}
\newunicodechar{≤}{\ensuremath{\leq}}
\newunicodechar{≥}{\ensuremath{\geq}}
\newunicodechar{−}{\ensuremath{-}}
\newunicodechar{×}{\ensuremath{\times}}
\newunicodechar{÷}{\ensuremath{\div}}
\newunicodechar{±}{\ensuremath{\pm}}
\newunicodechar{β}{\ensuremath{\beta}}
\newunicodechar{λ}{\ensuremath{\lambda}}
\newunicodechar{σ}{\ensuremath{\sigma}}

\usepackage{xcolor}
\definecolor{rnavy}{HTML}{1B365D}
\definecolor{rgrey}{HTML}{7A828C}
\definecolor{rline}{HTML}{BFC5CC}
\definecolor{raccent}{HTML}{C0392B}

\usepackage[a4paper,top=2.6cm,bottom=2.6cm,left=2.7cm,right=2.7cm,headheight=14pt]{geometry}
\linespread{1.45}
\setlength{\parskip}{0.35em}

% 标题层级：章（\section）、节（\subsection）、小节（\subsubsection，少用）
\ctexset{
  section={format=\sffamily\Large\bfseries\color{rnavy}\raggedright, name={,}, number=\arabic{section},
           aftername=\enspace, beforeskip=1.6em, afterskip=0.8em,
           numberformat={\color{white}\colorbox{rnavy}{\makebox[1.1em]{\arabic{section}}}}},
  subsection={format=\sffamily\large\bfseries\color{rnavy}\raggedright, beforeskip=1.2em, afterskip=0.5em},
  subsubsection={format=\sffamily\normalsize\bfseries\raggedright, beforeskip=0.9em, afterskip=0.3em},
}

% 页眉：左栏目、右日期，下方细色线；页脚：页码
\usepackage{fancyhdr}
\providecommand{\ReportHeadLeft}{策略研究报告}
\providecommand{\ReportHeadRight}{2026 年 11 月}
\pagestyle{fancy}
\fancyhf{}
\fancyhead[L]{\sffamily\small\color{rnavy}\ReportHeadLeft}
\fancyhead[R]{\sffamily\small\color{rnavy}\ReportHeadRight}
\renewcommand{\headrule}{{\color{rnavy}\hrule height 0.8pt}}
\fancyfoot[C]{\sffamily\small\thepage}
\fancypagestyle{plain}{\fancyhf{}\renewcommand{\headrulewidth}{0pt}\fancyfoot[C]{\sffamily\small\thepage}}

% 图表：标题在上，「图 1　标题」
\usepackage{graphicx,float}
\usepackage[font=small,labelfont={sf,bf,color=rnavy},textfont={sf,bf,color=rnavy},labelsep=colon,justification=raggedright,singlelinecheck=false,skip=5pt]{caption}
\DeclareCaptionLabelSeparator{zhcolon}{：}
\captionsetup{labelsep=zhcolon}
\captionsetup[figure]{name=图表,position=top}
\captionsetup[table]{name=图表,position=top}
\makeatletter\let\c@table\c@figure\makeatother
% 图表下方一行资料来源（唯一允许的小字）
\newcommand{\source}[1]{\par\vspace{2pt}{\raggedright\sffamily\footnotesize\color{rgrey}资料来源：#1\par}}

\usepackage{booktabs,tabularx,array,multirow}
\renewcommand{\arraystretch}{1.3}
\newcolumntype{L}{>{\raggedright\arraybackslash}X}
\newcolumntype{C}{>{\centering\arraybackslash}X}

\usepackage{amsmath,amssymb,bm}
\usepackage{enumitem}
\setlist{nosep,leftmargin=2em,itemsep=0.2em}

\usepackage{tikz}
\usetikzlibrary{arrows.meta,positioning,calc,fit,decorations.pathreplacing}
\tikzset{
  node/.style = {draw=rnavy, line width=0.6pt, rounded corners=1.5pt, align=flush center,
                 inner xsep=6pt, inner ysep=4pt, font=\sffamily\small},
  key/.style  = {node, draw=raccent, line width=0.8pt},
  soft/.style = {node, draw=rline},
  flow/.style = {-{Stealth[length=2mm]}, line width=0.6pt, draw=rnavy},
  note/.style = {font=\sffamily\small, text=rgrey, align=flush center},
}

% 首页研究要点：每条以一句加粗结论开头
\newenvironment{keypoints}{%
  \par\noindent{\sffamily\large\bfseries\color{rnavy}研究要点}\par\vspace{0.3em}
  {\color{rnavy}\hrule height 0.6pt}\vspace{0.6em}
  \begin{itemize}[leftmargin=1.2em,itemsep=0.6em,label={\color{rnavy}\rule[0.3ex]{0.45em}{0.45em}}]}
  {\end{itemize}\vspace{0.4em}{\color{rline}\hrule height 0.4pt}}
\newcommand{\kp}[1]{{\sffamily\bfseries\color{rnavy}#1}}

\usepackage[hidelinks]{hyperref}
\usepackage{bookmark}

% 封面
\makeatletter
\newcommand{\subtitle}[1]{\gdef\@subtitle{#1}}
\newcommand{\team}[1]{\gdef\@team{#1}}
\newcommand{\reportdate}[1]{\gdef\@reportdate{#1}}
\providecommand{\@subtitle}{}\providecommand{\@team}{}\providecommand{\@reportdate}{}
\newcommand{\makereporttitle}{%
  \thispagestyle{fancy}
  \vspace*{0.4em}
  {\raggedright\sffamily\color{rnavy}\fontsize{20pt}{27pt}\selectfont\bfseries\@title\par}
  \vspace{0.5em}
  {\raggedright\sffamily\large\color{rnavy}\@subtitle\par}
  \vspace{0.3em}
  {\raggedright\sffamily\small\color{rgrey}\@team\qquad\@reportdate\par}
  \vspace{1.2em}}
\makeatother

\title{因子赛道 AI 工具使用说明}
\subtitle{因子赛道决赛 · 提交材料}
\team{fx\hbox{战队}}
\reportdate{2026 年 11 月}

% 表格内部行距略收紧，只作用于本文件的表格
\newcommand{\tablesetup}{\small\linespread{1.05}\selectfont\renewcommand{\arraystretch}{1.18}}
\newcolumntype{P}[1]{>{\raggedright\arraybackslash}p{#1}}

\begin{document}
\makereporttitle

本说明披露因子赛道使用 AI 工具的情况。AI 编码代理编写了因子代码的主体、本地评测系统和实验脚本，批量构造并检验候选信号，起草了决赛提交材料。研究方向、评价标准、合规边界和提交作品由团队成员决定，比赛平台上的全部操作由团队成员完成。AI 的产出经第 4 节所列方式复核后才被采用，团队对全部提交内容负责。

\section{使用的工具与版本}

因子赛道主要使用了两种 AI 编码代理（图表 1）。主编码代理 Claude Code 全程使用；10 月 12 日至 18 日，批量编码代理按主编码代理撰写的任务说明与任务书生成候选因子，此后团队成员停用该代理，因子研究由主编码代理单独完成。

\begin{table}[!htbp]
\caption{AI 工具、版本与使用时段（2026 年）}
\centering
{\tablesetup
\begin{tabularx}{\textwidth}{@{}P{3.3cm}P{3.6cm}P{2.3cm}L@{}}
\toprule
工具 & 模型与版本 & 使用时段 & 运行方式与用途\\
\midrule
Claude Code（Anthropic），下称主编码代理 & claude-opus-5-5，客户端 2.1.293～2.1.321 & 10-09～10-31，11-10 起 & 在团队成员的本地计算机上运行，经授权操作团队自建服务器与工作站；编写代码、评测工具与实验脚本，撰写任务书并验收，起草决赛材料\\
Codex 命令行代理（OpenAI），下称批量编码代理 & 模型版本未留存记录 & 10-12～10-18 & 在团队自建服务器上运行；按任务书批量生成候选因子并提交本地评测\\
另一编码代理 & 名称与版本未留存记录 & 10-11 & 部署本地评测系统第一版，运行一批早期探索，结果未采用\\
\bottomrule
\end{tabularx}}
\source{团队整理}
\end{table}

\section{使用环节}

图表 2 按研究环节列出分工。AI 编码代理不接触比赛平台，平台返回的成绩与报错、探测脚本的运行结果和主办方的答复，均由团队成员取得后转交 AI 编码代理。

\begin{table}[!htbp]
\caption{各研究环节的分工}
\centering
{\tablesetup
\begin{tabularx}{\textwidth}{@{}P{2.4cm}LL@{}}
\toprule
环节 & AI 编码代理承担的工作 & 团队成员承担的工作\\
\midrule
评测口径 & 编写平台探测脚本与本地评测系统，核对官方标签与中性化口径 & 在平台运行探测脚本并取回结果，转达主办方规则\\
文献与方向 & 检索文献，整理待验证的方向 & 确定探索方向与优先次序\\
候选生成 & 撰写任务书与筛选流程，批量生成并评测候选 & 确定评价标准，判断探索是否过窄\\
组合与换手 & 构建基础组合，筛选增量信号，设计换手缓冲带 & 指出换手问题，追问树模型的过拟合与可解释性\\
合规检查 & 编写静态检查、截断重算与交付检查工具 & 追问未来信息来源，决定违规候选整组不提交\\
提交与选择 & 生成提交文件与说明，对比两种换手版本 & 在平台上传并回传成绩，选定私榜作品\\
记录与材料 & 维护得分记录，起草研究报告、展示材料、讲稿与本说明 & 审阅并提出修改意见，确定定稿\\
\bottomrule
\end{tabularx}}
\source{团队整理}
\end{table}

\section{关键提示词（转述）}

图表 3 按环节转述团队提出的 8 项关键要求及随后的结果。20 份任务书等较长的提示词由主编码代理按团队成员的要求起草，完整对话记录由团队留存。

\begin{table}[!htbp]
\caption{关键提示词的要求与随后的结果（2026 年，转述）}
\centering
{\tablesetup
\begin{tabularx}{\textwidth}{@{}P{1.6cm}P{1.0cm}P{6.3cm}L@{}}
\toprule
环节 & 日期 & 提示词要求 & 随后的结果\\
\midrule
评测口径 & 10-11 & 要求本地评测最大限度与平台保持同步 & 改为在模拟平台中原样运行提交代码\\
评价标准 & 10-14 & 要求不以榜单得分为目标，五项指标逐项改善 & 准入标准改为四项不变差、换手率不升\\
\midrule
候选生成 & 10-14 & 指出上千个候选多由少数几个信号改写而来，要求批量编码代理探索更多不同方向 & 统计出 1056 个因子中 942 个含同一反转分量，任务书改为寻找新的基础信号\\
 & 10-19 & 要求先扩展知识面，再做优化 & 按文献批量构造新信号，委托规模比由此进入最终因子\\
\midrule
合规 & 10-13 & 要求指出未来信息问题的具体位置，并判断修复后效果是否仍然成立 & 完成未来信息审计，停用约 240 个历史候选\\
 & 10-19 & 要求两个赛道的输出不得混用 & 否决把端到端模型预测混入因子\\
\midrule
换手控制 & 10-21 & 指出换手偏高，要求先批判现有做法再做探索，并在本地评测 & 分析评测的换手定义，设计换手缓冲带\\
提交选择 & 10-26 & 要求在表现相近的版本中选择换手不过低的一个 & 私榜两个作品分别取换手率约 0.28 与约 0.35 的版本\\
\bottomrule
\end{tabularx}}
\source{团队与 AI 编码代理的对话记录，团队整理}
\end{table}

\section{人工复核方式}

AI 的产出在进入提交前经过五个方面的复核。一是口径对齐：用平台处理后的因子值与官方标签复算平台自测，本地 RankIC 与平台相同，ICIR 与压力稳定性相差不超过 0.003。二是合规检查：静态检查拦截读取禁用数据表的代码，并要求风格暴露只用前一交易日；截断重算要求截止时刻之前的因子值逐位一致。三是交付检查：检查 Notebook 文件结构，并在模拟平台原样运行一个月、与本地值逐行比对，最终因子的 Notebook 176,000 行最大差异为 0。四是公榜事后检验，只与同一页面上的对照比较。五是由团队成员决定方向、评价标准、合规边界与私榜作品。复核也有局限：本地评测在多空组合夏普比率上与公榜的秩相关只有 0.59，排名靠前的候选之间由公榜同一页面的比较决定取舍；最终因子整体的截断重算没有留存记录，各分钟信号与竞价信号单独做过截断重算，结果全部一致。

第一个纠偏案例是标签口径。早期本地评测使用按公榜成绩反推的标签口径。团队成员发现本地与公榜的结论差别很大并提出质疑，AI 编码代理当时仍认为旧口径成立。随后 AI 编码代理编写平台探测脚本，由团队成员在平台运行，查明官方标签为截面后 30 分钟窗口的收盘价除以同窗口 VWAP 再减 1，逐行比对 100\% 吻合；旧口径与官方标签的截面秩相关只有 0.56。本地评测系统随即改用官方口径。

第二个案例是未来信息。10 月 12 日，批量编码代理报告的一组成绩突出的候选读取了评测用的风格暴露表；团队成员追问后核实，表中流通市值按当日收盘价计算，该组候选全部不提交，此后风格暴露只取前一交易日。10 月 13 日，团队成员追问另一批五项指标同时上升的候选，主编码代理当晚审计发现，日内最高价时段信号用整天数据计算时 RankIC 为 +0.019，只用截面之前的数据为 −0.035；按整天数据聚合的几类写法共涉及约 240 个历史候选，全部停用，截断重算由此加入筛选流程。

第三个案例是跨赛道混合。10 月 19 日，主编码代理把端到端模型的预测按 0.2 的权重混入因子做诊断，2023 年 RankIC 提高 0.0073，并提议采用；团队成员以两个赛道的输出不得混用为由否决，该组合从未进入交付。两个赛道的模型输出与因子值互相隔离，研究思路相互借鉴，第一个案例的标签核对即在端到端赛道完成。

\section{最终采用情况}

比赛期间共产生 1390 个因子编号，其中 1032 个由批量编码代理首次加入；本地评测任务按任务编号去重共 6600 个，其中包含同一候选的重复评测，不等于独立假设的个数；有公榜五项记录的提交 52 条；最终进入私榜的作品 2 个，即 1669 与 1670。

图表 4 列出最终因子各组成部分的来源。4 个同伴残差反转信号衡量个股相对走势最相似的一组股票的短期超涨超跌，是基础组合中分裂权重最大的一类输入，由团队成员在 AI 编码代理早期代码的基础上开发；其余代码由 AI 编码代理编写。

\begin{table}[!htbp]
\caption{最终因子各组成部分的来源}
\centering
{\tablesetup
\begin{tabularx}{\textwidth}{@{}P{3.4cm}LL@{}}
\toprule
组成部分 & 提出与实现 & 团队成员的作用\\
\midrule
同伴残差反转信号（4 个） & AI 编码代理编写早期版本，团队成员在此基础上开发；AI 编码代理接入基础组合并核对口径 & 开发 4 个版本，指定纳入基础组合\\
其余 23 个基础信号 & 两个编码代理提出与生成；批量编码代理的产出经主编码代理验收后改写 & 要求寻找新信息来源；确定合规边界\\
基础组合（浅层梯度提升树） & 主编码代理提出；训练脚本初版由批量编码代理编写 & 追问过拟合与可解释性；转达主办方答复\\
增量信号（15 项） & 主编码代理批量构造与筛选 & 确定准入标准\\
缓冲前的风格与行业中性化 & 主编码代理提出与实现 & 转达主办方关于风格暴露只用前一交易日的规则\\
换手缓冲带与缓冲强度 & 主编码代理分析评测的换手定义后设计，并对比两种换手版本 & 指出换手问题；提出私榜取舍方向，选定 1669 与 1670\\
09:30 截面输出 0 & 主编码代理自查后改正 & 质疑成绩可信度，取消相关任务\\
\bottomrule
\end{tabularx}}
\source{团队整理}
\end{table}

未被采用的 AI 产出同样有记录，主要有四类：含未来信息的候选，依据错误口径的早期探索，跨赛道混合提议，以及未通过公榜事后检验或无法在平台运行的改动，例如提高缓冲强度的版本和在平台上内存超限的备选版本。

最终提交的 1669 在私榜上 RankIC 为 0.0790、ICIR 为 1.559、换手率为 0.266。因子赛道公榜第二，私榜总榜第二。

\section{责任声明}

AI 编码代理不是本作品的作者。团队对本作品的代码、数据处理、实验结果、文字结论及其他全部提交内容承担责任。AI 编码代理不登录比赛平台，不向平台提交任何内容，平台上的全部操作由团队成员完成；AI 编码代理不修改已交付的提交文件，修复一律使用新的因子编号；不可逆的删除操作由团队成员执行。本说明与其他决赛材料由主编码代理及其调用的同模型子代理起草，经团队成员审阅修改后定稿。代码仓库提交历史、本地评测任务索引、实验日志、任务书、得分记录以及与 AI 编码代理的对话记录由团队留存，可应主办方要求提供。

\end{document}
